\clearpage\corrigesChapitre{11}{Décrire un mouvement}
\corrigesSouples
\begin{blocLivre}{cor-11-1}\begin{corrige}{1}{Un voyageur assis}
\begin{questions}
\item Elle est immobile par rapport à son siège, donc au train.
\item Elle se déplace avec le train par rapport au quai. Le référentiel doit accompagner toute description.
\end{questions}
\end{corrige}\end{blocLivre}
\begin{blocLivre}{cor-11-2}\begin{corrige}{2}{Choisir un point}
\begin{questions}
\item Il conserve le déplacement et la trajectoire du centre.
\item Il ne décrit ni la rotation du ballon ni sa déformation.
\end{questions}
\end{corrige}\end{blocLivre}
\begin{blocLivre}{cor-11-3}\begin{corrige}{3}{Une vitesse moyenne}
\begin{questions}
\item $\Delta x=8,0-2,0=\qty{6,0}{m}$.
\item $v_{\mathrm{moy}}=6,0/3,0=\qty{2,0}{m.s^{-1}}$, dans le sens positif.
\end{questions}
\end{corrige}\end{blocLivre}
\begin{blocLivre}{cor-11-4}\begin{corrige}{4}{Changer d’unité}
\begin{questions}
\item $72/3,6=\qty{20}{m.s^{-1}}$.
\item $d=v\Delta t=20\times0,50=\qty{10}{m}$.
\end{questions}
\end{corrige}\end{blocLivre}
\begin{blocLivre}{cor-11-5}\begin{corrige}{5}{Lire une chronophotographie}
\schema[corrige]{chronophoto}
\begin{questions}
\item La trajectoire est rectiligne ; les espacements croissent, donc la valeur de la vitesse augmente.
\item $v_1\simeq(0,80-0,20)/0,20=\qty{3,0}{m.s^{-1}}$, vers la droite.
\end{questions}
\end{corrige}\end{blocLivre}
\begin{blocLivre}{cor-11-6}\begin{corrige}{6}{Une vitesse négative}
\begin{questions}
\item $v=(2,0-5,0)/1,5=\qty{-2,0}{m.s^{-1}}$.
\item Valeur : \qty{2,0}{m.s^{-1}} ; sens des abscisses décroissantes. Le signe concerne l’orientation.
\end{questions}
\end{corrige}\end{blocLivre}
\begin{blocLivre}{cor-11-7}\begin{corrige}{7}{Échelle d’un vecteur}
\begin{questions}
\item $\ell=6,0/2,0=\qty{3,0}{cm}$.
\item Non, une longueur de vecteur représente une vitesse. Elle ne donne pas directement le déplacement de l’objet.
\end{questions}
\end{corrige}\end{blocLivre}
\begin{blocLivre}{cor-11-8}\begin{corrige}{8}{À valeur constante sur un cercle}
\schema[corrige]{mouvement-circulaire}
\begin{questions}
\item Non : sa direction change quand le point tourne.
\item Le vecteur vitesse est tangent à la trajectoire, orienté dans le sens du mouvement.
\end{questions}
\end{corrige}\end{blocLivre}
\begin{blocLivre}{cor-11-9}\begin{corrige}{9}{Un aller-retour}
\begin{questions}
\item Le déplacement total est nul, donc le vecteur vitesse moyenne est nul.
\item $40/10=\qty{4,0}{m.s^{-1}}$. Ce quotient mesure la rapidité moyenne sur le trajet, et non le déplacement vectoriel par durée.
\end{questions}
\end{corrige}\end{blocLivre}
\begin{blocLivre}{cor-11-10}\begin{corrige}{10}{Boucle de calcul}
\begin{questions}
\item 1 ; 3 ; 5 ; 7 \unit{m.s^{-1}}.
\item Chaque calcul utilise $i+1$ ; au dernier indice, aucune position suivante n’est disponible.
\end{questions}
\end{corrige}\end{blocLivre}
\begin{blocLivre}{cor-11-11}\begin{corrige}{11}{Une erreur de cadence}
\begin{questions}
\item Durée réelle : 0,020 s ; durée utilisée : 0,040 s, deux fois plus grande.
\item Pour un même déplacement, la vitesse calculée est divisée par deux.
\end{questions}
\end{corrige}\end{blocLivre}
\begin{blocLivre}{cor-11-12}\begin{corrige}{12}{Deux dimensions}
\demarche{Déterminer les composantes du vecteur, puis sa norme par Pythagore.}
\begin{questions}
\item $v_x=0,60$ et $v_y=\qty{0,80}{m.s^{-1}}$.
\item $v=\sqrt{0,60^2+0,80^2}=\qty{1,0}{m.s^{-1}}$, par Pythagore.
\end{questions}
\retiens{Les composantes d’une vitesse ne s’additionnent pas pour donner sa valeur.}
\end{corrige}\end{blocLivre}

